% !TEX root = ../main.tex
% Section: Appendix: Tightened Conditional Momentum Rates

\section{Tightened Conditional Momentum Rates}
\label{app:tightened-conditional-momentum}\label{app:momentum-details}

This appendix sharpens Lemma~\ref{lem:conditional}. It gives the exact
optimizer within the same Lyapunov estimates, the simpler parameter
choice used in the numerical figures, and the corresponding global
bootstrap with its energy prefactor.

\begin{proposition}[Tightened conditional energy decay]
\label{prop:tightened-conditional-momentum-pl}\label{prop:tightened}
Under the loss and position bounds of Lemma~\ref{lem:conditional} on
$[0,T]$, fix $\eta\in(0,1)$ and $\varepsilon>0$. With its constants
$\mathfrak K,\mathfrak L,\rho$, set
\begin{align}
 q&:=\varepsilon\sqrt{\mathfrak L/2},\label{eq:tightened-momentum-q-definition}\\
 r_{\varepsilon,\eta}
 &:=\min\{\varepsilon(1-\eta)\mathfrak K,\,
          2\gamma-\varepsilon(\rho+\gamma^2/(2\eta))\},\label{eq:tightened-momentum-dissipation-definition}\\
 \lambda_{\varepsilon,\eta}
 &:=r_{\varepsilon,\eta}/(1+q).\label{eq:tightened-momentum-rate-definition}
\end{align}
If $q<1$ and $r_{\varepsilon,\eta}>0$, the corrected energy
$\mathcal V_t:=\cE_t+\varepsilon\dot\Delta_t$ satisfies, for $0\le t\le T$,
\begin{align}
 (1-q)\cE_t&\le\mathcal V_t\le(1+q)\cE_t,\label{eq:tightened-momentum-lyapunov-equivalence}\\
 \mathcal V_t&\le\mathcal V_0e^{-\lambda_{\varepsilon,\eta}t},\label{eq:tightened-conditional-momentum-rate}\\
 \ell(\Svu_t)-\ell_\star
 &\le\cE_t\le\frac{\mathcal V_0}{1-q}e^{-\lambda_{\varepsilon,\eta}t}
 \le\frac{1+q}{1-q}\cE_0e^{-\lambda_{\varepsilon,\eta}t}.\label{eq:tightened-conditional-momentum-loss-rate}
\end{align}
If $\dot\Delta_0=0$, the first energy prefactor is $\cE_0/(1-q)$.
\end{proposition}

\begin{proof}
The force bounds \eqref{eq:conditional-force-bounds} give
$|\dot\Delta|\le\sqrt{2\mathfrak L\Delta\mathcal{K}}
\le\sqrt{\mathfrak L/2}\,\cE$, proving the equivalence.
For the derivative, use $\dot{\cE}=-2\gamma\mathcal{K}$,
$\ddot\Delta\le\rho\mathcal{K}-\mathcal N-\gamma\dot\Delta$, and Young's inequality:
\begin{equation}\label{eq:tightened-momentum-parameterized-young}
 \gamma|\dot\Delta|\le\gamma\sqrt{2\mathcal{K}\mathcal N}
 \le\eta\mathcal N+\frac{\gamma^2}{2\eta}\mathcal{K}.
\end{equation}
Since $\mathcal N\ge\mathfrak K\Delta$, this yields
\[
 \dot{\mathcal V}
 \le-\varepsilon(1-\eta)\mathcal N
     -[2\gamma-\varepsilon(\rho+\gamma^2/(2\eta))]\mathcal{K}
 \le-r_{\varepsilon,\eta}\cE
 \le-\lambda_{\varepsilon,\eta}\mathcal V.
\]
Integrate on $[0,T]$ and use the equivalence once more.
\end{proof}

\begin{proposition}[Optimal rate within the Lyapunov bounds]
\label{prop:optimal-tightened-conditional-momentum-rate}
Under Proposition~\ref{prop:tightened}, fix $\eta\in(0,1)$ and write
\begin{equation}\label{eq:tightened-momentum-A-B-definition}
 A_\eta:=(1-\eta)\mathfrak K,\qquad
 B_\eta:=\rho+\gamma^2/(2\eta).
\end{equation}
The unique maximizing correction at fixed $\eta$ is admissible and gives
\begin{align}
 \varepsilon_\eta^\star&=\frac{2\gamma}{A_\eta+B_\eta},\label{eq:tightened-momentum-optimal-epsilon}\\
 q_\eta^\star&:=\varepsilon_\eta^\star\sqrt{\mathfrak L/2}<1,\label{eq:tightened-momentum-optimal-q}\\
 \lambda_\eta&:=\frac{2\gamma A_\eta}
 {A_\eta+B_\eta+\gamma\sqrt{2\mathfrak L}}.\label{eq:tightened-momentum-fixed-eta-rate}
\end{align}
The unique maximizer over $\eta$ and its rate are
\begin{align}
 \eta_\star&=\frac{\gamma}
 {\gamma+\sqrt{\gamma^2+2\rho+2\gamma\sqrt{2\mathfrak L}}}
 \in(0,1/2),\label{eq:tightened-momentum-optimal-eta}\\
 \lambda_\star&:=\max_{0<\eta<1}\lambda_\eta
 =\frac{2\gamma\mathfrak K}
 {\mathfrak K+\frac12
 (\gamma+\sqrt{\gamma^2+2\rho+2\gamma\sqrt{2\mathfrak L}})^2}.\label{eq:tightened-momentum-optimal-rate}
\end{align}
\end{proposition}

\begin{proof}
At fixed $\eta$, the rate is the minimum of
$\varepsilon A_\eta/(1+\varepsilon\sqrt{\mathfrak L/2})$ and
$(2\gamma-\varepsilon B_\eta)/(1+\varepsilon\sqrt{\mathfrak L/2})$.
The first is strictly increasing and the second strictly decreasing,
so their intersection gives the unique optimum.
It is admissible because
$B_\eta>\rho+\gamma^2/2\ge\gamma\sqrt{2\rho}
\ge\gamma\sqrt{2\mathfrak L}$, hence $q_\eta^\star<1$;
both dissipation terms there equal $\varepsilon_\eta^\star A_\eta>0$.

For the remaining optimization, set
$\widetilde\rho:=\rho+\gamma\sqrt{2\mathfrak L}$ and divide numerator
and denominator of \eqref{eq:tightened-momentum-fixed-eta-rate} by $1-\eta$.
Maximizing the rate is equivalent to minimizing
$(\widetilde\rho+\gamma^2/(2\eta))/(1-\eta)$.
Its derivative has the sign of
$\widetilde\rho\eta^2+\gamma^2\eta-\gamma^2/2$, which increases strictly
and vanishes only at \eqref{eq:tightened-momentum-optimal-eta}.
The quotient there is
$\frac12(\gamma+\sqrt{\gamma^2+2\widetilde\rho})^2$, proving the formula.
\end{proof}

The numerical figures use the simpler admissible choice
\begin{equation}\label{eq:tightened-momentum-eta-approx}
 \eta_{\mathrm{approx}}
 :=\frac{\gamma}{\gamma+\sqrt{\gamma^2+2\rho}}\in(0,1/2).
\end{equation}
It minimizes $B_\eta/A_\eta$: differentiation gives
$\rho\eta^2+\gamma^2\eta-\gamma^2/2=0$.
Thus it optimizes the dissipation rate before accounting for the
equivalence factor $1+q_\eta^\star$. That factor is retained in the
following rigorous bound.

\begin{corollary}[Explicit rate used in the figures]
\label{cor:explicit-tightened-conditional-momentum-rate}\label{cor:tightened-choices}
Under Proposition~\ref{prop:tightened}, use
$\eta_{\mathrm{approx}}$ from \eqref{eq:tightened-momentum-eta-approx} and set
\begin{align}
 A_{\mathrm{approx}}&:=(1-\eta_{\mathrm{approx}})\mathfrak K,&
 B_{\mathrm{approx}}&:=\rho+\frac{\gamma^2}{2\eta_{\mathrm{approx}}},\label{eq:tightened-momentum-explicit-A-B}\\
 \varepsilon_{\mathrm{approx}}&:=\frac{2\gamma}{A_{\mathrm{approx}}+B_{\mathrm{approx}}},&
 q_{\mathrm{approx}}&:=\varepsilon_{\mathrm{approx}}\sqrt{\mathfrak L/2}.\label{eq:tightened-momentum-explicit-epsilon-q}
\end{align}
These parameters are admissible. With
$\mathcal V_0=\cE_0+\varepsilon_{\mathrm{approx}}\dot\Delta_0$,
\begin{align}
 \lambda_{\mathrm{approx}}
 &:=\frac{\varepsilon_{\mathrm{approx}}A_{\mathrm{approx}}}
         {1+q_{\mathrm{approx}}},\label{eq:tightened-momentum-explicit-rate}\\
 \ell(\Svu_t)-\ell_\star
 &\le\cE_t\le\frac{\mathcal V_0}{1-q_{\mathrm{approx}}}
       e^{-\lambda_{\mathrm{approx}}t},\qquad 0\le t\le T.\label{eq:tightened-momentum-explicit-loss-rate}
\end{align}
If $\dot\Delta_0=0$, the prefactor is $\cE_0/(1-q_{\mathrm{approx}})$.
\end{corollary}

\begin{proof}
The fixed-$\eta$ optimum in
Proposition~\ref{prop:optimal-tightened-conditional-momentum-rate}
is admissible for every $\eta\in(0,1)$.
Apply Proposition~\ref{prop:tightened} with $\eta=\eta_{\mathrm{approx}}$.
\end{proof}

The approximate choice concerns only optimization: its bound is exact,
and $\lambda_{\mathrm{approx}}\le\lambda_\star$. Neither exponent is
claimed to be the optimal rate of the underlying training dynamics.

\begin{corollary}[Global bootstrap with the correct prefactor]
\label{cor:tightened-bootstrap}
Under the loss and initialization hypotheses of Theorem~\ref{thm:momentum},
choose a tube $a_s:=a_0e^{-2s}$, $b_s:=b_0e^{2s}$ and an admissible
$(\varepsilon,\eta)$ for its constants. Put
\[
 \mathfrak C_0:=\frac{\cE_0+\varepsilon\dot\Delta_0}{(1-q)\cE_0}.
\]
If $\cE_0>0$ and
\begin{equation}\label{eq:tightened-bootstrap}
 \frac2{\lambda_{\varepsilon,\eta}}
 \sqrt{\frac{2\mathfrak C_0\cE_0}{a_s}}<s,
\end{equation}
then the tube and \eqref{eq:tightened-conditional-momentum-loss-rate} hold globally.
A sufficient condition using only initial position bounds and total
energy replaces $\mathfrak C_0$ by $(1+q)/(1-q)$.
\end{corollary}
\begin{proof}
Within the initial connected positive-position chart,
$\int_0^t\norm D_{\op}\le\int_0^t\sqrt{2\cE_r/a_s}\,\dd r$.
Proposition~\ref{prop:tightened} bounds this by the left side of
\eqref{eq:tightened-bootstrap}. The strict first-exit proof of
Theorem~\ref{thm:momentum} then applies unchanged. Initial positivity and
the Lyapunov equivalence imply $1\le \mathfrak C_0\le(1+q)/(1-q)$.
\end{proof}

\paragraph{The role of the prefactor.}
A faster conditional exponent must be accompanied by its own prefactor
when testing the global deformation budget.
The sharper state-dependent prefactor uses the initial loss derivative
$\dot\Delta_0=\langle E_0,\Sigma^{qu}_0+\Sigma^{vp}_0\rangle$
as well as total energy. In particular, zero predictor speed gives
$\mathfrak C_0=1/(1-q)$. The parameter-only prefactor uses no extra initial derivative.
Maximizing the conditional exponent alone need not minimize the global
budget: the same choice of parameters also changes $q$ and hence the
prefactor. The closed-form optimizer should therefore be tested in
\eqref{eq:tightened-bootstrap}, not substituted there with an unchanged
prefactor.

\section{Large-Position-Scale Certificate Asymptotics}
\label{app:large-position-scale-certificate}\label{app:large-position-scale}
\label{subsec:large-position-scale-certificate}

Increasing position variance can open a damping window without reducing
the initial loss gap. We quantify that tradeoff and the best rates
certified by the same bootstrap.
For the two-scale white initialization of
Corollary~\ref{cor:two-scale-white-window}, the initial mechanical excess
energy is
\[
\cE_0
=
\ell(0)-\ell_\star+\frac{m\sigma_{qp}^2}{2}.
\]
This quantity is fixed when \(\sigma_{qp}\) is fixed and
\(\sigma_{vu}\to+\infty\).  It is useful to name the bootstrap gap
\[
\mathfrak G_s(\gamma):=\mathfrak S_s(\gamma)-s,
\qquad
s_{\mathrm{cert}}(\gamma)
:=
\inf\{s>0:\mathfrak G_s(\gamma)<0\}.
\]
Whenever this set is nonempty, define the corresponding boundary rate by
\[
\lambda_{\mathrm{cert}}(\gamma)
:=
\lambda_{s_{\mathrm{cert}}(\gamma)}(\gamma),
\]
using continuity in \(s\).  The strict bootstrap condition need not hold at
the boundary itself: the theorem applies at interior budgets
\(s>s_{\mathrm{cert}}\) with \(\mathfrak G_s(\gamma)<0\), whose rates approach
this boundary value.  Thus \(\lambda_{\mathrm{cert}}\) is a limiting supremal
certificate, rather than a rate asserted at a possibly nonadmissible boundary
point.
Define also the largest damping for which the bootstrap window is nonempty:
\[
\gamma^{\mathrm{cert}}_{\max}(\sigma_{vu})
:=
\sup\{\gamma>0:\text{ there exists }s>0
\text{ such that }\mathfrak G_s(\gamma)<0\}.
\]

The assumption \(\cE_0>0\) below excludes only the trivial zero-excess
case, in which both the predictor loss gap and kinetic energy vanish and the
trajectory is stationary.

\paragraph{Explicit damping interval and feasibility.}
At a fixed deformation budget, admissibility reduces to a quadratic
inequality, making both damping boundaries explicit.
For \(\cE_0>0\) and fixed deformation budget \(s>0\), write
\begin{equation}\label{eq:window}
 c_s:=\frac{s\kappa a_s^{3/2}}{4\sqrt{3\cE_0/2}}.
\end{equation}
The theorem's strict certificate is exactly
\(2\gamma^2-c_s\gamma+\rho_s<0\), or
\begin{equation}\label{eq:window-roots}
 c_s^2>8\rho_s,\qquad
 \frac{c_s-\sqrt{c_s^2-8\rho_s}}4<\gamma<
 \frac{c_s+\sqrt{c_s^2-8\rho_s}}4.
\end{equation}
These are sufficient boundaries, not instability thresholds.


\paragraph{Optimizing feasibility.}
To make this interval nonempty as easily as possible, maximize
\begin{equation}\label{eq:budget}
 \frac{c_s^2}{8\rho_s}
 =\frac{s^2\kappa^2a_0^3e^{-6s}}
 {192\cE_0(4b_0Le^{2s}+2\sqrt{2L\cE_0})}.
\end{equation}
Its unique maximizer $s_{\rm feas}$ belongs to $(1/4,1/3)$ and satisfies
\begin{equation}\label{eq:best-budget}
 \frac1{s_{\rm feas}}
 =3+\frac{4b_0Le^{2s_{\rm feas}}}
 {4b_0Le^{2s_{\rm feas}}+2\sqrt{2L\cE_0}}.
\end{equation}
For a proof, write the denominator term as $Ae^{2s}+B$, with $A,B>0$.
The logarithmic derivative is $2/s-6-2Ae^{2s}/(Ae^{2s}+B)$, and its
derivative is $-2/s^2-4ABe^{2s}/(Ae^{2s}+B)^2<0$.
It changes sign between $1/4$ and $1/3$. The resulting scalar equation
is explicit and uniquely solvable; no trajectory has to be computed.
This budget optimizes \emph{nonemptiness}, not the upper damping endpoint
or the decay exponent. Those are different optimization questions.


\paragraph{Asymptotic certificate bounds.}
The large-position limit distinguishes fixed-damping behavior from the
largest damping that the certificate can accommodate.
\begin{proposition}[Large-position-scale certificate asymptotics]
\label{prop:large-position-scale-certificate}\label{prop:large-position}
For the two-scale white initialization above, fix \(\sigma_{qp}\ge0\), assume
\(\cE_0>0\), and
let \(\sigma_{vu}\to+\infty\).  For each fixed \(\gamma>0\) and fixed
\(s>0\),
\[
\mathfrak S_s(\gamma)
=
\frac{16Le^{5s}}{\kappa\gamma\,\sigma_{vu}}
\sqrt{\frac{3\cE_0}{2}}
+O(\sigma_{vu}^{-3}).
\]
The expansion is uniform for \(s\) in bounded sets.  Consequently,
\[
s_{\mathrm{cert}}(\gamma)
=
\frac{16L}{\kappa\gamma\,\sigma_{vu}}
\sqrt{\frac{3\cE_0}{2}}
+o(\sigma_{vu}^{-1}),
\]
and
\[
\lambda_{\mathrm{cert}}(\gamma)
=
\frac{\kappa\gamma}{4L}
\;-\;
\frac{16}{\sigma_{vu}}\sqrt{\frac{3\cE_0}{2}}
+o(\sigma_{vu}^{-1}).
\]
Moreover, the upper damping boundary satisfies
\[
\gamma^{\mathrm{cert}}_{\max}(\sigma_{vu})
=
\frac{\kappa\sigma_{vu}^3}
{24e\sqrt{3\cE_0/2}}
\bigl(1+o(1)\bigr).
\]
\end{proposition}

\begin{proof}
For fixed $\gamma>0$, the exact expression \eqref{eq:certificate}
can be expanded directly. Since
\[
\rho_s=4L\sigma_{vu}^2e^{2s}
+2\sqrt{2L\cE_0},
\]
we have
\[
\lambda_s(\gamma)
=
\frac{\kappa\gamma}{4Le^{4s}}
+O(\sigma_{vu}^{-2}).
\]
Substituting this into the definition of \(\mathfrak S_s(\gamma)\), with
\(a_s=\sigma_{vu}^2e^{-2s}\), gives the first expansion.
Write
\[
c_0:=\frac{16L}{\kappa\gamma}
\sqrt{\frac{3\cE_0}{2}}.
\]
The expansion is uniform near \(s=0\).  Thus every
\(s=c/\sigma_{vu}\), with \(c>0\) fixed, satisfies
\[
\mathfrak G_s(\gamma)
=\frac{c_0-c}{\sigma_{vu}}+o(\sigma_{vu}^{-1}),
\]
uniformly when \(c\) ranges over a compact set.  Hence budgets whose leading
constant is below \(c_0\) are inadmissible, whereas those whose leading
constant is strictly above \(c_0\) are admissible for all sufficiently large
\(\sigma_{vu}\).  This proves the displayed expansion for
\(s_{\mathrm{cert}}(\gamma)\).  Evaluating
\[
\lambda_s(\gamma)
=
\frac{\kappa\gamma}{4L}e^{-4s}
+O(\sigma_{vu}^{-2})
\]
at \(s=s_{\mathrm{cert}}(\gamma)\) gives the expansion for
\(\lambda_{\mathrm{cert}}(\gamma)\).

It remains to identify the high-damping boundary.  First observe the exact
upper bound
\[
\lambda_s(\gamma)
=\frac{a_s\kappa\gamma}{\rho_s+2\gamma^2}
\le\frac{\kappa\sigma_{vu}^2e^{-2s}}{2\gamma}.
\]
Consequently every admissible pair satisfies
\[
\mathfrak S_s(\gamma)
\ge
\frac{8\gamma e^{3s}}{\kappa\sigma_{vu}^3}
\sqrt{\frac{3\cE_0}{2}},
\]
and hence
\[
\gamma
<
\frac{\kappa\sigma_{vu}^3}{8\sqrt{3\cE_0/2}}\,s e^{-3s}
\le
\frac{\kappa\sigma_{vu}^3}{24e\sqrt{3\cE_0/2}}.
\]
This proves the required asymptotic upper bound uniformly over all budgets,
not only over bounded \(s\).

For the matching lower bound, take \(\gamma\) of order
\(\sigma_{vu}^3\) and keep \(s\) in a compact subset of \((0,+\infty)\).
Then $2\gamma^2$ dominates $\rho_s$ in \eqref{eq:certificate}, and
\[
\lambda_s(\gamma)
=
\frac{\kappa\sigma_{vu}^2e^{-2s}}{2\gamma}
\bigl(1+o(1)\bigr).
\]
Consequently,
\[
\mathfrak S_s(\gamma)
=
\frac{8\gamma e^{3s}}{\kappa\sigma_{vu}^3}
\sqrt{\frac{3\cE_0}{2}}
\bigl(1+o(1)\bigr),
\]
so the condition \(\mathfrak S_s(\gamma)<s\) is asymptotically equivalent to
\[
\gamma
<
\frac{\kappa\sigma_{vu}^3}{8\sqrt{3\cE_0/2}}\,
s e^{-3s}.
\]
The right-hand side is maximized at \(s=1/3\).  Choosing this budget and any
fixed multiplicative factor below the maximum gives an admissible pair for all
sufficiently large \(\sigma_{vu}\), proving the matching lower bound and the
claimed asymptotic formula.
\end{proof}

Proposition~\ref{prop:large-position-scale-certificate} has two complementary
messages.  At fixed damping, increasing the initial position scale makes the
certified ellipticity deformation vanish and the certified rate converges to
\(\kappa\gamma/(4L)\).  At fixed \(\sigma_{vu}\), however, this
particular bootstrap theorem cannot certify arbitrarily large damping: the
admissible region in \((s,\gamma)\) eventually disappears, with the upper
edge scaling like \(\sigma_{vu}^3\).  This is a limitation of the certificate,
not a claim that the underlying damped dynamics fail in the large-damping
regime.

\begin{figure}[t]
  \centering
  \def\mathdefault#1{#1}
  \subimport{figures/}{momentum_bootstrap_lambda_multi_sigma.pgf}
\caption{
Limiting bootstrap rate at the admissibility boundary as a function of $\gamma$ for several initial position
scales $\sigma_{vu}$.
Solid curves are obtained by grid-searching the bootstrap certificate
for the boundary value \(\lambda_{s_{\mathrm{cert}}}(\gamma)\); rates at
strictly admissible nearby budgets approach these curves.  The dotted black curve is
the large-\(\sigma_{vu}\) limiting rate \(\kappa\gamma/(4L)\).  The
plot uses \((d_u, d_v) = (3, 2)\), \(\kappa=L=1\), \(\sigma_{qp}=1\), and
\(\Delta_0=\ell(0)-\ell_\star=1\).
For scales whose upper damping boundary lies within the scanned range, markers show numerical and asymptotic approximations to the maximal admissible damping and its corresponding boundary rate.
}
\label{fig:momentum-bootstrap-lambda-multi-sigma}\label{fig:momentum-boundary}
\end{figure}

\FloatBarrier

%%% Local Variables:
%%% mode: LaTeX
%%% TeX-master: "../main"
%%% End:
